\chapter{Evaluation}\label{chap:evaluation}

This chapter evaluates the new Mitos against the design requirements of \Cref{sec:design-reqs}. Requirements \ref{req:3}, \ref{req:4} and \ref{req:10} (interface) were evaluated qualitatively against the original Mitos C API in \Cref{sec:design-reqs}. Requirements \ref{req:5}--\ref{req:9} are evaluated as functional checks (\Cref{sec:eval-functional}). Requirements \ref{req:1}, \ref{req:2}, \ref{req:7} are evaluated as data-fidelity comparisons (\Cref{sec:eval-fidelity}); \ref{req:16} and \ref{req:17} as performance comparisons (\Cref{sec:eval-performance}). Requirements \ref{req:11}--\ref{req:15} (microarchitecture independence) are argued by design and supported empirically by the cross-platform replication of the same evaluation (\Cref{sec:eval-archindep}).

\section{Methodology}\label{sec:eval-methodology}

\paragraph{Platforms.} Two systems are used: an Intel Skylake-SP system (Xeon Bronze 3106) and an Intel Ice Lake-SP system (Xeon Platinum 8360Y), both exercising the Intel PEBS path. Both Mitos variants are expected to work on both systems. Ice Lake was originally identified as a platform on which original Mitos did not run, which was a big part of the initial motivation for the work, but a later investigation found that it does in fact run there; the cause of the earlier failure is unknown and was not reproduced in time for this evaluation. The Ice Lake column therefore reports a full comparison rather than a one-sided measurement, and \ref{req:13} (PEBS on Ice Lake) is verified by demonstrating that new Mitos produces the expected output on Ice Lake. An AMD system was originally planned in order to exercise both IBS backends (\ref{req:12}); two attempts on different machines were made, but neither completed. One machine did not grant the necessary \texttt{kernel.perf\_event\_paranoid} permissions, and the Caliper installation on the other ran into Spack-related build issues that could not be resolved in time. The empirical IBS column is therefore absent from the result tables; the IBS claim rests on the design argument in \Cref{sec:eval-archindep} alone, and the consequences for \ref{req:12} are listed in \Cref{sec:eval-threats}.

\paragraph{Workload.} A naive triple-nested matrix multiplication, in serial, OpenMP-parallel, and MPI-parallel variants is profiled under each Mitos variant. The matmul kernel is shipped as an example in original Mitos's source tree, and was simply adapted for the new implementation. The computation stays exactly the same. Compilation flags are \texttt{-O3 -g -std=c++17}; \texttt{-g} is required by \ref{req:7}. The matrix dimension is fixed at $N = 1024$ for the principal results; OpenMP uses $4$ threads, MPI uses $4$ ranks.

\paragraph{Procedure.} For each cell of the experimental grid (platform $\times$ implementation variant $\times$ Mitos variant), $R = 20$ repetitions are run, preceded by a repetition as warm-up which is discarded. The machine is otherwise idle. For each cell the median across surviving repetitions is reported.

\paragraph{Configuration parity.} For comparisons regarding \ref{req:1} and \ref{req:2}, the two Mitos variants are configured to sample the same period. The kernel-level sample source is therefore held constant by construction, and any observed difference in the sample stream is attributable to the new code path (libpfm4, Caliper services, and \texttt{cali\_to\_csv.py}).

\section{Functional checks (R5--R9)}\label{sec:eval-functional}

\Cref{tab:eval-functional} reports the pass/fail status of each functional requirement on each platform, for both Mitos variants.

\ref{req:8} (sample callback) was not tested as part of this. This is because the implementation differs. The new Mitos creates a new Caliper Attribute for the collected value, which is not output by default, as to not deviate from the output format of the original Mitos.
For this implementation, the callback functionality works as described previously.

\ref{req:9} (MemAxes format) was not tested. The new implementation replicates the output format generated by the original Mitos, including a \texttt{hwloc}-provided hardware.xml and copying of source files to a subdirectory. Compatibility is therefore assumed but not verified.

\begin{table}[h]
    \centering
    \small
    \begin{tabular}{@{}clcccc@{}}
        \textbf{Req} & \textbf{What is checked}                & \multicolumn{2}{c}{\textbf{Skylake}} & \multicolumn{2}{c}{\textbf{Ice Lake}}                         \\
                     &                                         & orig                                 & new                                   & orig & new            \\
        \hline
        \ref{req:5}  & every OpenMP worker contributes samples & yes                                  & no                                    & yes  & no             \\
        \ref{req:6}  & every MPI rank contributes samples      & yes                                  & no                                    & yes  & no             \\
        \ref{req:7}  & sample carries non-empty source/line    & yes                                  & yes (seq, omp)                        & yes  & yes (seq, omp) \\
        \ref{req:9}  & MemAxes loads the output without errors & yes                                  & unsure                                & yes  & unsure         \\
    \end{tabular}
    \label{tab:eval-functional}
\end{table}

The serial path runs end-to-end for both Mitos variants on both platforms; the differences there are quantitative (sample counts, function-name overlap) rather than functional, and are taken up in \Cref{sec:eval-fidelity}. The OpenMP path collects samples on the new Mitos but only from the master thread; the MPI path performed similiarly at one point in time but currently collects no samples at all.

\section{Data (R1, R2, R7)}\label{sec:eval-fidelity}

\Cref{tab:eval-fidelity} reports, per cell, the median sample count across surviving repetitions for both Mitos variants, the cross-tool top-10 function Jaccard, and the source-resolution rate.

\begin{table}[h]
    \centering
    \small
    \begin{tabular}{@{}llrrcrr@{}}
        \textbf{Platform} & \textbf{Variant} & \multicolumn{2}{c}{\textbf{n\_samples}} & \multicolumn{2}{c}{\textbf{resolved \%}}                \\
                          &                  & orig                                    & new                                      & orig & new   \\
        \hline
        Skylake           & seq              & 5262                                    & 3146                                     & 99.2 & 100.0 \\
        Skylake           & omp              & 3156                                    & 841                                      & 99.7 & 100.0 \\
        Skylake           & mpi              & 10613                                   & 0                                        & 13.8 & ---   \\
        Ice Lake          & seq              & 5293                                    & 3301                                     & 99.9 & 100.0 \\
        Ice Lake          & omp              & 3238                                    & 845                                      & 99.9 & 100.0 \\
        Ice Lake          & mpi              & 17242                                   & 0                                        & 26.3 & ---   \\
    \end{tabular}
    \caption{Data-fidelity metrics per platform and parallelism variant; medians across 20 repetitions. ``---'' in the new-Mitos resolved-\% column indicates that no samples were collected.}
    \label{tab:eval-fidelity}
\end{table}

The sample-count comparison gives the cleanest read on data quantity (R2). In the serial cells the new Mitos collects 60--62\,\% of the original's samples (3146 of 5262 on Skylake, 3301 of 5293 on Ice Lake) at the same configured period --- a substantial undercount. In the OpenMP cells the gap is much larger (${\sim}27\,\%$), but most of this is assumed to be a cascade effect from only the master thread contributing samples, so this is effectively one-thread sampling compared against four-thread sampling. The MPI cells are an outright failure.

Caliper consistently reports more distinct \texttt{(file,\,line)} tuples (3--10 vs.\ 1--2), which is consistent with preserved-vs-collapsed inlines. Whether the underlying samples agree on the hot function is therefore not directly answered.

Source-resolution (R7) is at 100\,\% on every new-Mitos cell where samples were collected, which matches or slightly exceeds the original on serial and OpenMP. I am unsure wether the poor MPI baseline at 13.8\,\% on Skylake and 26.3\,\% on Ice Lake is due to an issue with my tests.

\section{Performance (R16, R17)}\label{sec:eval-performance}

\Cref{tab:eval-performance} reports median wall-clock runtime, for both Mitos variants.

\begin{table}[h]
    \centering
    \small
    \begin{tabular}{@{}llrr@{}}
        \textbf{Platform} & \textbf{Variant} & \textbf{Runtime [s] (orig/new)} \\
        \hline
        Skylake           & seq              & 7.71 / 7.78                     \\
        Skylake           & omp              & 7.82 / 2.83                     \\
        Skylake           & mpi              & 4.09 / 3.92                     \\
        Ice Lake          & seq              & 4.67 / 4.98                     \\
        Ice Lake          & omp              & 4.73 / 1.65                     \\
        Ice Lake          & mpi              & 2.83 / 2.76                     \\
    \end{tabular}
    \caption{Performance metrics per platform and parallelism variant; medians across 20 repetitions. The OpenMP runtime advantage for new Mitos make it seem like the original Mitos was not parallelized, but this is unverified.}
    \label{tab:eval-performance}
\end{table}

For the serial workload, new-Mitos runtime is within 1\,\% of the original on Skylake (7.78 vs.\ 7.71\,s) and within 7\,\% on Ice Lake (4.98 vs.\ 4.67\,s); \ref{req:17} is met on serial. For OpenMP, new Mitos's runtime is sharply lower than the original (2.83 vs.\ 7.82 on Skylake; 1.65 vs.\ 4.73 on Ice Lake), but the gap is likely due to improper parallelization for the original Mitos. This is not a real performance improvement. For MPI, runtimes are comparable (within ${\sim}5\,\%$ on both platforms), but new Mitos's runtime is over a workload that collects no samples, so the comparison says nothing about real sampling overhead.

\section{Microarchitecture independence (R11--R15)}\label{sec:eval-archindep}

The microarchitecture-independence requirements are not amenable to direct measurement on a finite set of platforms: the universal claim ``works on architectures I have not tested'' cannot be empirically discharged. I therefore argue that the design meets the requirements (aside from the other reported issues), and substantiate the argument with the cross-platform replication reported in the preceding sections.

\paragraph{Argument by design.} New Mitos contains no vendor-specific PMU encoding in its own source. The PMU event string is interpreted by libpfm4 (\Cref{sec:libpfm4}), which maintains compiled-in descriptor tables for Intel, AMD, ARM, and several other vendors and translates the same human-readable event string to the appropriate \texttt{perf\_event\_attr} encoding per microarchitecture. Adding support for a new microarchitecture in the same vendor family typically requires only a libpfm4 update; adding a new sampling backend (e.g.\ ARM SPE) requires a small configuration entry in the Caliper libpfm service. Original Mitos's hard-coded \texttt{perf\_event\_attr} construction (\Cref{sec:mitos}) precludes this.

\paragraph{Empirical support.} The same source compiles, links, and produces output of the documented shape on both test platforms (\Cref{tab:eval-functional}). The Skylake and Ice Lake runs share the PEBS code path and differ only in the libpfm4-emitted encoding. \ref{req:13} (PEBS on Ice Lake) is verified by demonstrating that new Mitos produces the expected output on Ice Lake. \ref{req:12} (Intel PEBS and AMD IBS support) is met for PEBS empirically and for IBS by design only, as the AMD system test was not completed (\Cref{sec:eval-threats}). Both tested microarchitectures are Intel, so the empirical evidence speaks to within-vendor portability; the cross-vendor claim rests on the design argument alone.

\section{Threats to validity}\label{sec:eval-threats}

\begin{itemize}
    \item AMD/IBS path not tested. \ref{req:12} (Intel PEBS and AMD IBS support) is verified for PEBS only. Two attempts to test new Mitos on AMD systems failed. The IBS half of \ref{req:12} therefore rests on the design argument in \Cref{sec:eval-archindep} alone.
    \item MPI and OpenMP tested separately. No MPI+OpenMP test was run, so the claim that new Mitos supports both in the same run is unverified. The MPI and OpenMP paths were tested separately, and both failed to reach feature parity with the original Mitos; the combined path is therefore not expected to work well either.
    \item Architecture independence is argued, not proved. Only two microarchitectures, both Intel, are tested. The claim about future hardware rests on the libpfm4 abstraction layer (\Cref{sec:libpfm4}); the empirical work substantiates it only insofar as the same source compiles and runs without modification on both test platforms.
    \item Cross-tool Top-10 Jaccard is unreliable. The cross-tool Jaccard is 0 in every cell, which is implausible given matching workloads; this very likely reflects a name-normalisation mismatch in the analyzer rather than disjoint sampling. \ref{req:1} (data quality) therefore lacks a strong quantitative measure in this evaluation.
\end{itemize}